\documentclass{ximera}

\input{../../../preamble.tex}


\definecolor{fvdnavy}{HTML}{071D43}
\definecolor{fvdcyan}{HTML}{20BFC6}
\definecolor{fvdcyanlight}{HTML}{EFFCFB}
\definecolor{fvdmagenta}{HTML}{F10D69}
\definecolor{fvdmagentalight}{HTML}{FFF1F7}
\definecolor{fvdtext}{HTML}{163044}





%=================================================
% Pedagogische omgevingen
% PDF: tcolorbox
% HTML: learning-box
%=================================================

\newenvironment{voorkennis}
{%
    \ifdefined\HCode
        \HCode{%
<section class="learning-box learning-box--prior">
<div class="learning-box__header">
<span class="learning-box__icon" aria-hidden="true"></span>
<span class="learning-box__title">Voorkennis</span>
</div>
<div class="learning-box__content">
        }%
        \begin{itemize}
    \else
       \begin{tcolorbox}[
    enhanced,
    breakable,
    colback=fvdcyanlight,
    colframe=fvdcyan,
    coltext=fvdtext,
    boxrule=0.5pt,
    leftrule=4pt,
    arc=4mm,
    outer arc=4mm,
    left=7mm,
    right=6mm,
    top=5mm,
    bottom=5mm,
    before skip=6mm,
    after skip=6mm,
    title=Voorkennis,
    coltitle=fvdcyan,
    fonttitle=\bfseries\Large,
    attach title to upper={\par\smallskip},
    boxed title style={
        empty,
        left=0mm,
        right=0mm,
        top=0mm,
        bottom=0mm
    },
    overlay={
        \draw[
            fvdcyan,
            line width=0.5pt
        ]
        ([xshift=42mm,yshift=-13mm]frame.north west)
        --
        ([xshift=-7mm,yshift=-13mm]frame.north east);
    }
]
        \begin{itemize}
    \fi
}
{%
    \end{itemize}
    \ifdefined\HCode
        \HCode{%
</div>
</section>
        }%
    \else
        \end{tcolorbox}
    \fi
}


\newenvironment{leerdoelen}
{%
    \ifdefined\HCode
        \HCode{%
<section class="learning-box learning-box--goals">
<div class="learning-box__header">
<span class="learning-box__icon" aria-hidden="true"></span>
<span class="learning-box__title">Leerdoelen</span>
</div>
<div class="learning-box__content">
        }%
        \begin{itemize}
    \else
\begin{tcolorbox}[
    enhanced,
    breakable,
    colback=fvdmagentalight,
    colframe=fvdmagenta,
    coltext=fvdtext,
    boxrule=0.5pt,
    leftrule=4pt,
    arc=4mm,
    outer arc=4mm,
    left=7mm,
    right=6mm,
    top=5mm,
    bottom=5mm,
    before skip=6mm,
    after skip=6mm,
    title=Leerdoelen,
    coltitle=fvdmagenta,
    fonttitle=\bfseries\Large,
    attach title to upper={\par\smallskip},
    boxed title style={
        empty,
        left=0mm,
        right=0mm,
        top=0mm,
        bottom=0mm
    },
    overlay={
        \draw[
            fvdmagenta,
            line width=0.5pt
        ]
        ([xshift=42mm,yshift=-13mm]frame.north west)
        --
        ([xshift=-7mm,yshift=-13mm]frame.north east);
    }
]
        \begin{itemize}
    \fi
}
{%
    \end{itemize}
    \ifdefined\HCode
        \HCode{%
</div>
</section>
        }%
    \else
        \end{tcolorbox}
    \fi
}


\addPrintStyle{../../..}

\title{Vallen en werpen}
\author{Ingmar Herreman}

\begin{abstract}
Waarom valt een muntstuk sneller dan een los blad papier, maar niet
sneller dan datzelfde papier wanneer luchtweerstand geen rol speelt?
Waarom is de snelheid van een omhoog geworpen bal op het hoogste punt
nul, terwijl zijn versnelling daar niet nul is?

In dit hoofdstuk passen we de algemene kinematica uit het vorige
hoofdstuk toe op bewegingen onder invloed van de zwaartekracht.
We vertrekken van waarnemingen en maken vervolgens een fysisch model:
de vrije val. Daarbij maken we scherp onderscheid tussen een werkelijke
valbeweging en het geïdealiseerde model waarin andere krachten dan de
zwaartekracht worden verwaarloosd.

We analyseren vrije val en verticale worp met positie-, snelheids- en
versnellingsfuncties. We verbinden functievoorschriften met grafieken,
afgeleiden, de oppervlaktemethode en bepaalde integralen. Bij elk
probleem besteden we bijzondere aandacht aan het gekozen
referentiestelsel en de tekenconventie.

Zo gebruiken we dezelfde wiskundige structuur om een voorwerp dat
wordt losgelaten, omhoog geworpen of naar beneden geworpen te
beschrijven. Tegelijk leren we kritisch beoordelen wanneer het
vrije-valmodel bruikbaar is en wanneer de werkelijkheid een rijker
model vereist.
\end{abstract}

\begin{document}

% FVD_CHAPTER_DATA_START
% Automatisch gegenereerd uit index.tex.
% Niet handmatig aanpassen.
\fvdchapterdata{1}{Beweging en kracht}{3}
% FVD_CHAPTER_DATA_END

\maketitle


% ============================================================
% VOORKENNIS EN LEERDOELEN
% ============================================================

\begin{voorkennis}
  \item Je kan een rechtlijnige beweging beschrijven met $x(t)$, $v(t)$ en $a(t)$.
  \item Je kan snelheid bepalen als afgeleide van de plaatsfunctie.
  \item Je kan versnelling bepalen als afgeleide van de snelheidsfunctie.
  \item Je kan een EVRB analyseren met
        $v=v_0+at$ en
        $x=x_0+v_0t+\frac12at^2$.
  \item Je kan verplaatsing bepalen als georiënteerde oppervlakte onder een $v(t)$-grafiek.
  \item Je kan snelheidsverschil bepalen als georiënteerde oppervlakte onder een $a(t)$-grafiek.
  \item Je kan bepaalde integralen gebruiken om verplaatsing en snelheidsverschil te bepalen.
  \item Je kan een oorsprong en positieve richting kiezen en consequent met tekens werken.
\end{voorkennis}

\begin{leerdoelen}
  \item Je kan valbeweging en vrije val van elkaar onderscheiden.
  \item Je kan uitleggen onder welke voorwaarden het vrije-valmodel bruikbaar is.
  \item Je kan de valversnelling $\vec g$ fysisch interpreteren en het teken van $g_y$ bepalen uit een gekozen assenstelsel.
  \item Je kan verklaren waarom de massa in het ideale vrije-valmodel niet voorkomt in de bewegingsvergelijkingen.
  \item Je kan een vrije val zonder beginsnelheid analyseren met $y(t)$, $v_y(t)$ en $a_y(t)$.
  \item Je kan een verticale worp omhoog of omlaag analyseren met $y(t)$, $v_y(t)$ en $a_y(t)$.
  \item Je kan de grafieken van positie, snelheid en versnelling bij vrije val en verticale worp schetsen, analyseren en kwantificeren.
  \item Je kan het hoogste punt van een verticale worp correct interpreteren: $v_y=0$ maar $a_y\neq0$.
  \item Je kan verplaatsingen bij valbewegingen bepalen met de oppervlaktemethode en met bepaalde integralen.
  \item Je kan een valsituatie mathematiseren door een schets, referentiestelsel, gegevens en geschikt model te kiezen.
  \item Je kan resultaten kritisch beoordelen en de beperkingen van het model zonder luchtweerstand aangeven.
  \item Je kan gegevens of videobeelden van een valbeweging gebruiken om het model van een EVRB te onderzoeken.
\end{leerdoelen}


% ============================================================
% 1 OPENING
% ============================================================

\FVDsection{Wat bepaalt hoe een voorwerp valt?}

Laat van dezelfde hoogte tegelijk een tennisbal en een pingpongbal
vallen. Herhaal het experiment met een muntstuk en een los blad papier.

Wat verwacht je?

In het dagelijkse leven lijkt het vanzelfsprekend dat sommige
voorwerpen sneller vallen dan andere. Een muntstuk bereikt de grond
duidelijk eerder dan een los blad papier. Daaruit besluiten dat
``zwaardere voorwerpen sneller vallen'' is echter te snel.

Voer een tweede proef uit. Leg het blad papier op het muntstuk en laat
beide samen vallen. Het papier kan nu vrijwel tegelijk met het
muntstuk de grond bereiken.

Er is dus meer aan de hand dan alleen massa.

\begin{oefening}[Voorspel voor je rekent]{\oefB\ESS}

Beantwoord de vragen vóór je verdere theorie leest.

\begin{question}
Welke van twee bollen met dezelfde vorm maar verschillende massa
verwacht je in vacuüm eerst op de grond?
\end{question}

\begin{question}
Waarom kan een los blad papier in lucht veel trager vallen dan een
muntstuk?
\end{question}

\begin{question}
Welke variabele zou je in een experiment veranderen om het effect van
luchtweerstand zichtbaar te maken?
\end{question}

\end{oefening}


% ============================================================
% 2 VALBEWEGING EN VRIJE VAL
% ============================================================

\FVDsection{Valbeweging en vrije val zijn niet hetzelfde}

Een \textbf{valbeweging} is een beweging waarbij een voorwerp onder
invloed van de zwaartekracht naar beneden beweegt. In een werkelijke
situatie kunnen daarnaast andere krachten optreden, vooral
luchtweerstand.

Een blad papier dat door de lucht naar beneden dwarrelt, voert dus een
valbeweging uit, maar zijn versnelling is niet voortdurend gelijk aan
de valversnelling.

Voor veel problemen gebruiken we een eenvoudiger model.

\begin{definitie}[Vrije val]

Bij een \textbf{vrije val} beschrijven we een voorwerp nadat het is
losgelaten of geworpen terwijl we aannemen dat alleen de zwaartekracht
de beweging beïnvloedt.

Dicht bij het aardoppervlak behandelen we de valversnelling daarbij
als constant.

\end{definitie}

In dit hoofdstuk betekent ``vrije val'' dus niet noodzakelijk
``vallen vanuit rust''. Ook een bal die verticaal omhoog wordt geworpen,
bevindt zich na het loslaten in vrije val wanneer we luchtweerstand
verwaarlozen.

\begin{opmerking}

Het woord \emph{vrij} verwijst naar het model waarin alleen de
zwaartekracht werkzaam is. Een voorwerp kan dus omhoog bewegen en
toch in vrije val zijn.

\end{opmerking}


\FVDsubsection{Werkelijkheid en model}

We gebruiken het schema

\[
\boxed{
\text{werkelijke beweging}
\longrightarrow
\text{aannames}
\longrightarrow
\text{vrije-valmodel}
}
\]

met als belangrijke aannames:

\begin{itemize}
  \item de luchtweerstand is verwaarloosbaar;
  \item de beweging vindt plaats dicht genoeg bij het aardoppervlak;
  \item de verandering van $g$ over de beschouwde hoogte is verwaarloosbaar;
  \item het voorwerp kan voor de beweging als puntmassa worden behandeld.
\end{itemize}

Onder die voorwaarden is de verticale versnelling constant en is de
beweging een EVRB.


\begin{oefening}[Vrije val of niet?]{\oefB\ESS}

Beoordeel of het vrije-valmodel waarschijnlijk bruikbaar is. Motiveer.

\begin{question}
Een kleine stalen kogel valt $1,0\,\mathrm{m}$ in een klaslokaal.
\end{question}

\begin{question}
Een open valscherm daalt naar de grond.
\end{question}

\begin{question}
Een verfrommeld blad papier valt van een tafel.
\end{question}

\begin{question}
Een veertje valt door een vacuümbuis.
\end{question}

\begin{question}
Een regendruppel valt vanuit een wolk tot op de grond.
\end{question}

\end{oefening}


\begin{oefening}[Een model bekritiseren]{\oefU\ESS}

Een leerling gebruikt voor een valschermspringer die van
$2000\,\mathrm{m}$ hoogte springt gedurende de volledige beweging het
model

\[
a=-9,81\,\mathrm{m/s^2}.
\]

\begin{question}
Welke belangrijke aanname maakt de leerling?
\end{question}

\begin{question}
Waarom wordt die aanname na verloop van tijd steeds minder realistisch?
\end{question}

\begin{question}
Wat verwacht je dat er met de versnelling gebeurt wanneer de
luchtweerstand belangrijk wordt?
\end{question}

\end{oefening}


% ============================================================
% 3 VALVERSNELLING
% ============================================================

\FVDsection{De valversnelling}

Dicht bij het aardoppervlak is de grootte van de valversnelling
ongeveer

\[
\boxed{
g=9,81\,\mathrm{m/s^2}.
}
\]

De vector $\vec g$ is naar het middelpunt van de aarde gericht.

De waarde $9,81\,\mathrm{m/s^2}$ is de \textbf{grootte} van de
valversnelling en is daarom positief. Het teken van de component van
de versnelling hangt af van de gekozen positieve richting.

\begin{opmerking}

Maak onderscheid tussen

\[
g=9,81\,\mathrm{m/s^2}
\]

als grootte van de valversnelling en bijvoorbeeld

\[
a_y=-g=-9,81\,\mathrm{m/s^2}
\]

wanneer de positieve $y$-richting omhoog is gekozen.

\end{opmerking}


\FVDsubsection{Waarom valt massa weg uit het model?}

In het ideale model is de zwaartekracht op een voorwerp dicht bij het
aardoppervlak

\[
F_z=mg.
\]

Volgens de tweede wet van Newton:

\[
F_{\mathrm{res}}=ma.
\]

Wanneer de zwaartekracht de enige kracht is:

\[
ma=mg.
\]

Voor $m\neq0$ volgt:

\[
\boxed{a=g}
\]

in grootte.

De massa valt dus uit de vergelijking. Dat verklaart waarom twee
voorwerpen met verschillende massa in hetzelfde zwaarteveld dezelfde
valversnelling hebben wanneer andere krachten verwaarloosbaar zijn.

\begin{opmerking}

Dit is een verbinding met de dynamica: de kinematica beschrijft
\emph{hoe} het voorwerp beweegt; de tweede wet van Newton verklaart
\emph{waarom} de versnelling in dit model onafhankelijk is van de
massa.

\end{opmerking}


\begin{oefening}[Zwaar of licht?]{\oefU\ESS}

Twee massieve bollen hebben dezelfde vorm en afmetingen. Bol B heeft
een tweemaal zo grote massa als bol A. Ze worden in vacuüm gelijktijdig
van dezelfde hoogte losgelaten.

\begin{question}
Vergelijk de zwaartekrachten op beide bollen.
\end{question}

\begin{question}
Vergelijk hun versnellingen.
\end{question}

\begin{question}
Bereiken ze de grond op hetzelfde tijdstip? Verklaar met de tweede
wet van Newton.
\end{question}

\end{oefening}


% ============================================================
% 4 TEKENCONVENTIE
% ============================================================

\FVDsection{Eerst kiezen, dan rekenen}

Veel fouten bij valbewegingen zijn geen rekenfouten maar
\textbf{tekenfouten}. Daarom kiezen we vóór elke berekening:

\begin{enumerate}
  \item een oorsprong;
  \item een positieve richting;
  \item een symbool voor de verticale positie.
\end{enumerate}

We gebruiken meestal de $y$-as.


\FVDsubsection{Omhoog positief}

Kiezen we omhoog als positieve richting, dan geldt altijd:

\[
\boxed{a_y=-g.}
\]

Een voorwerp dat omhoog beweegt heeft

\[
v_y>0,
\]

een voorwerp dat naar beneden beweegt heeft

\[
v_y<0.
\]


\FVDsubsection{Naar beneden positief}

Kiezen we naar beneden als positieve richting, dan:

\[
\boxed{a_y=+g.}
\]

Een vallend voorwerp heeft dan een positieve snelheid.

Beide keuzes zijn correct. Alleen de tekens verschillen.

\begin{eigenschap}[Consequentie is belangrijker dan de keuze]

Een fysisch antwoord mag niet afhangen van de gekozen positieve
richting. De tussenresultaten kunnen andere tekens hebben, maar
meetbare grootheden zoals valtijd en de grootte van de impactsnelheid
moeten dezelfde zijn.

\end{eigenschap}


\begin{oefening}[Tekencheck]{\oefB\ESS}

Kies de positieve $y$-richting omhoog.

Geef het teken van $v_y$ en $a_y$ voor:

\begin{question}
een bal die omhoog beweegt na te zijn geworpen;
\end{question}

\begin{question}
dezelfde bal op het hoogste punt;
\end{question}

\begin{question}
dezelfde bal tijdens het dalen.
\end{question}

\end{oefening}


\begin{oefening}[Twee assenstelsels, één fysica]{\oefU}

Een steen wordt vanuit rust van een hoogte $h$ losgelaten.

Analyseer dezelfde beweging tweemaal:

\begin{question}
met omhoog als positieve richting;
\end{question}

\begin{question}
met naar beneden als positieve richting.
\end{question}

Vergelijk de tekens van positie, verplaatsing, snelheid en versnelling.
Leg uit waarom de berekende valtijd dezelfde moet zijn.

\end{oefening}


% ============================================================
% 5 VRIJE VAL UIT RUST
% ============================================================

\FVDsection{Vrije val zonder beginsnelheid}

We kiezen in deze sectie de positieve $y$-richting omhoog.

Een voorwerp wordt op tijdstip $t=0$ losgelaten vanaf hoogte $y_0$.

Dan:

\[
v_{y,0}=0
\]

en

\[
a_y=-g.
\]

De algemene EVRB-vergelijkingen uit het vorige hoofdstuk worden:

\[
v_y(t)=v_{y,0}+a_yt
\]

en

\[
y(t)=y_0+v_{y,0}t+\frac12a_yt^2.
\]

Dus:

\[
\boxed{
v_y(t)=-gt
}
\]

en

\[
\boxed{
y(t)=y_0-\frac12gt^2.
}
\]

De verplaatsing na tijd $t$ is:

\[
\boxed{
\Delta y=-\frac12gt^2.
}
\]

De grootte van de snelheid is:

\[
\boxed{
|v_y|=gt.
}
\]


\FVDsubsection{De drie grafieken}

Bij vrije val uit rust met omhoog positief:

\begin{itemize}
  \item $a_y(t)=-g$ is een horizontale rechte onder de tijdas;
  \item $v_y(t)=-gt$ is een dalende rechte door de oorsprong;
  \item $y(t)=y_0-\frac12gt^2$ is een dalende parabool.
\end{itemize}

De drie grafieken zijn niet onafhankelijk:

\[
\boxed{
v_y(t)=y'(t),
\qquad
a_y(t)=v_y'(t)=y''(t).
}
\]


\begin{voorbeeld}[Een steen van een klif]

Een steen wordt van een klif losgelaten en raakt
$1,80\,\mathrm{s}$ later het water. Verwaarloos luchtweerstand.

We kiezen omhoog positief en nemen het wateroppervlak als $y=0$.

De hoogte van de klif volgt uit:

\[
0=y_0-\frac12g(1,80)^2.
\]

Dus:

\[
y_0
=
\frac12\cdot9,81\cdot(1,80)^2
\approx15,9\,\mathrm{m}.
\]

De snelheid bij het raken van het water:

\[
v_y(1,80)
=
-9,81\cdot1,80
\approx-17,7\,\mathrm{m/s}.
\]

Het minteken betekent dat de snelheid naar beneden gericht is.
De grootte van de impactsnelheid is:

\[
\boxed{|v_y|\approx17,7\,\mathrm{m/s}.}
\]

\end{voorbeeld}


\begin{oefening}[Vrije val uit rust]{\oefB\ESS}

Een kleine metalen bol wordt van $20,0\,\mathrm{m}$ hoogte losgelaten.
Verwaarloos luchtweerstand.

\begin{question}
Maak een schets en kies een assenstelsel.
\end{question}

\begin{question}
Stel $y(t)$ en $v_y(t)$ op.
\end{question}

\begin{question}
Bepaal de valtijd.
\end{question}

\begin{question}
Bepaal de snelheid vlak vóór het raken van de grond.
\end{question}

\begin{question}
Interpreteer het teken van de snelheid.
\end{question}

\end{oefening}


% ============================================================
% 6 OPPERVLAKTE EN INTEGRAAL
% ============================================================

\FVDsection{Vrije val in het $v(t)$- en $a(t)$-diagram}

Hoofdstuk 2 gaf ons twee krachtige verbanden:

\[
\Delta y
=
\int_{t_1}^{t_2}v_y(t)\,dt
\]

en

\[
\Delta v_y
=
\int_{t_1}^{t_2}a_y(t)\,dt.
\]

Bij vrije val worden deze relaties bijzonder overzichtelijk.


\FVDsubsection{Van versnelling naar snelheid}

Met omhoog positief is:

\[
a_y=-g.
\]

De $a(t)$-grafiek vormt over een interval $\Delta t$ een rechthoek
onder de tijdas.

De georiënteerde oppervlakte is:

\[
\Delta v_y=-g\Delta t.
\]

Analytisch:

\[
\Delta v_y
=
\int_{t_1}^{t_2}-g\,dt
=
-g(t_2-t_1).
\]


\FVDsubsection{Van snelheid naar verplaatsing}

Bij een val uit rust is:

\[
v_y(t)=-gt.
\]

De $v(t)$-grafiek vormt een driehoek onder de tijdas. Over $[0,t]$:

\[
\Delta y
=
-\frac12\cdot t\cdot gt
=
-\frac12gt^2.
\]

Met een integraal:

\[
\Delta y
=
\int_0^t -g\tau\,d\tau
=
-\frac12gt^2.
\]

De bewegingsvergelijking, de oppervlaktemethode en de integraal geven
dus exact dezelfde fysische relatie.


\begin{eigenschap}[Drie talen, één beweging]

Voor vrije val kunnen we dezelfde beweging analyseren met:

\[
\boxed{
\text{bewegingsvergelijkingen}
\quad\longleftrightarrow\quad
\text{grafiek en oppervlakte}
\quad\longleftrightarrow\quad
\text{bepaalde integraal}.
}
\]

Een sterke analyse kan tussen deze voorstellingen schakelen.

\end{eigenschap}


\begin{oefening}[Dezelfde val op drie manieren]{\oefU\ESS}

Een voorwerp wordt vanuit rust losgelaten en valt gedurende
$2,00\,\mathrm{s}$. Kies omhoog positief.

\begin{question}
Bepaal de snelheidsverandering met de oppervlakte onder de
$a(t)$-grafiek.
\end{question}

\begin{question}
Bepaal de verplaatsing met de oppervlakte onder de $v(t)$-grafiek.
\end{question}

\begin{question}
Controleer beide resultaten met bepaalde integralen.
\end{question}

\begin{question}
Controleer de verplaatsing ten slotte met de EVRB-formule.
\end{question}

\end{oefening}


% ============================================================
% 7 VERTICALE WORP
% ============================================================

\FVDsection{Verticale worp}

Een voorwerp wordt verticaal omhoog of omlaag geworpen. Zodra het
voorwerp de hand verlaat en we luchtweerstand verwaarlozen, werkt in
ons model alleen de zwaartekracht.

De beweging is opnieuw een EVRB, maar nu is de beginsnelheid in het
algemeen niet nul.

Met omhoog positief:

\[
\boxed{a_y=-g}
\]

en dus:

\[
\boxed{
v_y(t)=v_{y,0}-gt
}
\]

en

\[
\boxed{
y(t)=y_0+v_{y,0}t-\frac12gt^2.
}
\]

Deze twee functies beschrijven zowel een worp omhoog als een worp
omlaag. Het teken van $v_{y,0}$ bepaalt de aanvankelijke
bewegingsrichting.


\FVDsubsection{Worp omhoog}

Bij een worp omhoog:

\[
v_{y,0}>0.
\]

Tijdens het stijgen:

\[
v_y>0,
\qquad
a_y<0.
\]

Snelheid en versnelling hebben tegengestelde tekens, dus de grootte
van de snelheid neemt af.


\FVDsubsection{Het hoogste punt}

Op het hoogste punt geldt:

\[
\boxed{v_y=0.}
\]

Maar:

\[
\boxed{a_y=-g.}
\]

De versnelling is daar dus niet nul.

\begin{opmerking}

``Het voorwerp staat op het hoogste punt even stil'' mag niet worden
verward met ``er is geen versnelling''.

De snelheid is op dat ene ogenblik nul. De snelheid verandert echter
nog steeds met $-9,81\,\mathrm{m/s}$ per seconde.

\end{opmerking}


\FVDsubsection{Tijdens het dalen}

Na het hoogste punt:

\[
v_y<0,
\qquad
a_y<0.
\]

Snelheid en versnelling hebben nu hetzelfde teken. Daarom neemt de
grootte van de snelheid toe.


\begin{oefening}[Stijgen, hoogste punt, dalen]{\oefB\ESS}

Een bal wordt verticaal omhoog geworpen. Omhoog is positief.

Vul voor elke fase het teken van $v_y$ en $a_y$ in en vermeld of
$|v_y|$ toeneemt, afneemt of ogenblikkelijk nul is.

\[
\begin{array}{c|c|c|c}
\text{fase} & v_y & a_y & |v_y|\\ \hline
\text{stijgen} & & &\\
\text{hoogste punt} & & &\\
\text{dalen} & & &
\end{array}
\]

\end{oefening}


\begin{oefening}[Misconcepten]{\oefU\ESS}

Beoordeel elke uitspraak. Verbeter ze indien nodig.

\begin{question}
Tijdens het stijgen is de versnelling positief omdat de bal omhoog beweegt.
\end{question}

\begin{question}
Op het hoogste punt zijn snelheid en versnelling beide nul.
\end{question}

\begin{question}
Tijdens het dalen is de versnelling negatief wanneer omhoog positief gekozen is.
\end{question}

\begin{question}
Een negatieve snelheid betekent dat het voorwerp vertraagt.
\end{question}

\end{oefening}


% ============================================================
% 8 HOOGSTE PUNT
% ============================================================

\FVDsection{Het hoogste punt kwantificeren}

Een bal wordt vanaf hoogte $y_0$ omhoog geworpen met beginsnelheid
$v_0>0$. We kiezen omhoog positief.


\FVDsubsection{Tijd tot het hoogste punt}

Op het hoogste punt:

\[
v_y=0.
\]

Uit

\[
v_y(t)=v_0-gt
\]

volgt:

\[
0=v_0-gt_{\max}.
\]

Dus:

\[
\boxed{
t_{\max}=\frac{v_0}{g}.
}
\]


\FVDsubsection{Hoogtewinst}

We kunnen de hoogtewinst bepalen door $t_{\max}$ in de plaatsfunctie
in te vullen. Een tweede methode gebruikt het $v(t)$-diagram.

Tijdens het stijgen vormt het gebied onder de $v(t)$-grafiek een
driehoek:

\[
\Delta y_{\max}
=
\frac12 v_0t_{\max}.
\]

Met

\[
t_{\max}=\frac{v_0}{g}
\]

krijgen we:

\[
\boxed{
\Delta y_{\max}
=
\frac{v_0^2}{2g}.
}
\]

Dus:

\[
\boxed{
y_{\max}
=
y_0+\frac{v_0^2}{2g}.
}
\]

Deze relatie maakt onmiddellijk zichtbaar dat de hoogtewinst
evenredig is met het \textbf{kwadraat} van de beginsnelheid.


\begin{voorbeeld}[Een bal omhoog werpen]

Een bal verlaat een hand op

\[
y_0=1,50\,\mathrm{m}
\]

met

\[
v_0=13,0\,\mathrm{m/s}.
\]

De tijd tot het hoogste punt is:

\[
t_{\max}
=
\frac{13,0}{9,81}
\approx1,33\,\mathrm{s}.
\]

De hoogtewinst:

\[
\Delta y_{\max}
=
\frac{13,0^2}{2\cdot9,81}
\approx8,61\,\mathrm{m}.
\]

Dus:

\[
y_{\max}
\approx
1,50+8,61
=
10,1\,\mathrm{m}.
\]

\end{voorbeeld}


\begin{oefening}[Dubbele beginsnelheid]{\oefU\ESS}

Bal A en bal B worden vanaf dezelfde hoogte verticaal omhoog geworpen.
De beginsnelheid van B is tweemaal die van A:

\[
v_{0,B}=2v_{0,A}.
\]

\begin{question}
Vergelijk de tijden tot het hoogste punt.
\end{question}

\begin{question}
Vergelijk de hoogtewinsten.
\end{question}

\begin{question}
Verklaar waarom een dubbele beginsnelheid niet leidt tot een dubbele
hoogtewinst.
\end{question}

\end{oefening}


\begin{oefening}[Algemeen verband afleiden]{\oefG}

Leid zonder de plaatsfunctie te gebruiken de formule

\[
\Delta y_{\max}=\frac{v_0^2}{2g}
\]

af uit het $v(t)$-diagram en de oppervlaktemethode.

Geef daarna aan hoe dezelfde formule met een bepaalde integraal kan
worden verkregen.

\end{oefening}


% ============================================================
% 9 SYMMETRIE
% ============================================================

\FVDsection{Omhoog en omlaag: een bijzondere symmetrie}

Wanneer een voorwerp omhoog wordt geworpen en later terugkeert naar
\textbf{dezelfde hoogte}, ontstaat in het ideale model een belangrijke
symmetrie.

Stel dat het voorwerp vanaf $y_0$ wordt geworpen met snelheid $v_0$.

Op dezelfde hoogte tijdens het dalen geldt:

\[
v_y=-v_0.
\]

De richting is omgekeerd, maar de grootte is dezelfde:

\[
|v_y|=v_0.
\]

Ook de tijd omhoog is gelijk aan de tijd terug naar dezelfde hoogte:

\[
t_{\mathrm{terug}}=2t_{\max}.
\]

\begin{opmerking}

Deze symmetrie geldt alleen wanneer we dezelfde begin- en eindhoogte
vergelijken en het model zonder luchtweerstand gebruiken.

Wanneer het voorwerp op een andere hoogte landt, mag je de totale
bewegingstijd niet zomaar verdubbelen.

\end{opmerking}


\begin{oefening}[Wat is symmetrisch?]{\oefU\ESS}

Een bal wordt vanaf een balkon omhoog geworpen en landt uiteindelijk
op de straat, die lager ligt dan het balkon.

Beoordeel:

\begin{question}
Is de tijd van loslaten tot het hoogste punt gelijk aan de tijd van
het hoogste punt tot de straat?
\end{question}

\begin{question}
Is de snelheid bij het opnieuw passeren van de balkonhoogte gelijk
aan de beginsnelheid?
\end{question}

\begin{question}
Is de \emph{grootte} van die snelheid gelijk aan de beginsnelheid?
\end{question}

Motiveer telkens.

\end{oefening}


% ============================================================
% 10 GRAFIEKEN
% ============================================================

\FVDsection{De drie grafieken analyseren}

Voor een verticale worp omhoog met omhoog positief:

\[
y(t)=y_0+v_0t-\frac12gt^2,
\]

\[
v_y(t)=v_0-gt,
\]

\[
a_y(t)=-g.
\]


\FVDsubsection{De $y(t)$-grafiek}

De plaatsfunctie is een neerwaarts geopende parabool.

\begin{itemize}
  \item De beginwaarde is $y_0$.
  \item Tijdens het stijgen neemt $y$ toe.
  \item Op het hoogste punt heeft $y(t)$ een maximum.
  \item De raaklijn is daar horizontaal omdat $v_y=0$.
  \item Daarna neemt $y$ af.
\end{itemize}


\FVDsubsection{De $v(t)$-grafiek}

De snelheidsfunctie is een dalende rechte met richtingscoëfficiënt

\[
-g.
\]

\begin{itemize}
  \item $v_y>0$ tijdens het stijgen;
  \item $v_y=0$ op het hoogste punt;
  \item $v_y<0$ tijdens het dalen.
\end{itemize}


\FVDsubsection{De $a(t)$-grafiek}

De versnelling is voortdurend:

\[
a_y=-g.
\]

De grafiek is dus een horizontale rechte onder de tijdas.


\begin{oefening}[Van één grafiek naar drie]{\oefU\ESS}

Een $v(t)$-grafiek van een verticale worp is een rechte die start bij

\[
v(0)=18\,\mathrm{m/s}
\]

en een richtingscoëfficiënt

\[
-9,81\,\mathrm{m/s^2}
\]

heeft.

\begin{question}
Bepaal het tijdstip van het hoogste punt.
\end{question}

\begin{question}
Schets de bijbehorende $a(t)$-grafiek.
\end{question}

\begin{question}
Schets kwalitatief de $y(t)$-grafiek.
\end{question}

\begin{question}
Duid op alle grafieken hetzelfde hoogste punt aan.
\end{question}

\begin{question}
Leg uit welke grafische eigenschap op elke grafiek met dat punt
overeenkomt.
\end{question}

\end{oefening}


% ============================================================
% 11 VERTICAAL OMLAAG
% ============================================================

\FVDsection{Verticaal naar beneden werpen}

Een verticale worp hoeft niet omhoog gericht te zijn.

Kiezen we omhoog positief en werpen we een voorwerp naar beneden, dan:

\[
v_{y,0}<0
\]

en nog steeds:

\[
a_y=-g.
\]

Dus:

\[
v_y(t)=v_{y,0}-gt.
\]

Snelheid en versnelling zijn vanaf het begin beide negatief. De grootte
van de snelheid neemt daarom toe.


\begin{oefening}[Steen naar beneden]{\oefB\ESS}

Vanop een platform $30,0\,\mathrm{m}$ boven de grond wordt een steen
verticaal naar beneden geworpen met een beginsnelheid van
$4,00\,\mathrm{m/s}$.

Kies omhoog positief en de grond als oorsprong.

\begin{question}
Noteer $y_0$, $v_{y,0}$ en $a_y$ met correct teken.
\end{question}

\begin{question}
Stel $y(t)$ en $v_y(t)$ op.
\end{question}

\begin{question}
Bepaal wanneer de steen de grond bereikt.
\end{question}

\begin{question}
Bepaal de impactsnelheid.
\end{question}

\end{oefening}


% ============================================================
% 12 PROBLEEMOPLOSSEN
% ============================================================

\FVDsection{Valproblemen systematisch oplossen}

Bij vrije val en verticale worp zijn de formules meestal niet het
moeilijkste onderdeel. De grootste uitdaging is de situatie correct
mathematiseren.

Gebruik daarom deze aanpak:

\begin{enumerate}
  \item \textbf{Schets de situatie.}
  \item \textbf{Kies oorsprong en positieve richting.}
  \item \textbf{Noteer de beginvoorwaarden:} $y_0$, $v_{y,0}$ en $a_y$.
  \item \textbf{Noteer de gezochte grootheid.}
  \item \textbf{Vertaal bijzondere gebeurtenissen:}
        hoogste punt $\Rightarrow v_y=0$,
        grond $\Rightarrow y=y_{\mathrm{grond}}$, enzovoort.
  \item \textbf{Kies een methode:} bewegingsvergelijking, grafiek,
        oppervlaktemethode of integraal.
  \item \textbf{Los op en selecteer fysisch zinvolle oplossingen.}
  \item \textbf{Controleer teken, eenheid en grootteorde.}
  \item \textbf{Interpreteer het antwoord in woorden.}
\end{enumerate}


\begin{voorbeeld}[Een tennisbal vanuit een raam]

Vanuit een raam op $8,00\,\mathrm{m}$ boven de grond wordt een
tennisbal verticaal omhoog geworpen. De bal raakt na
$2,00\,\mathrm{s}$ de grond. Welke beginsnelheid had de bal?

We kiezen de grond als $y=0$ en omhoog positief.

\[
y_0=8,00\,\mathrm{m},
\qquad
y(2,00)=0,
\qquad
a_y=-9,81\,\mathrm{m/s^2}.
\]

Uit

\[
y(t)=y_0+v_0t-\frac12gt^2
\]

volgt:

\[
0
=
8,00+2,00v_0
-\frac12\cdot9,81\cdot(2,00)^2.
\]

Dus:

\[
2,00v_0=11,62
\]

en:

\[
\boxed{
v_0\approx5,81\,\mathrm{m/s}.
}
\]

Het positieve teken bevestigt dat de bal aanvankelijk omhoog werd
geworpen.

\end{voorbeeld}


\begin{oefening}[Van raam naar straat]{\oefU\ESS}

Een bal wordt vanaf $12,0\,\mathrm{m}$ boven de grond verticaal omhoog
geworpen met $v_0=8,00\,\mathrm{m/s}$.

\begin{question}
Maak een schets en stel $y(t)$ en $v_y(t)$ op.
\end{question}

\begin{question}
Bepaal de maximale hoogte boven de grond.
\end{question}

\begin{question}
Bepaal wanneer de bal de grond raakt.
\end{question}

\begin{question}
Bij het oplossen van $y(t)=0$ ontstaan algebraïsch twee oplossingen.
Bespreek welke oplossingen fysisch betekenisvol zijn.
\end{question}

\begin{question}
Bepaal de snelheid bij impact.
\end{question}

\end{oefening}


\begin{oefening}[De baksteen]{\oefU\ESS}

Een bouwvakker werpt een baksteen verticaal omhoog vanaf
$1,50\,\mathrm{m}$ boven de grond. Een metselaar wil de steen op een
stelling op $4,00\,\mathrm{m}$ hoogte kunnen opvangen.

\begin{question}
Welke minimale beginsnelheid is nodig om die hoogte te bereiken?
\end{question}

\begin{question}
Welke bijzondere voorwaarde geldt bij de minimale beginsnelheid op
$4,00\,\mathrm{m}$ hoogte?
\end{question}

\begin{question}
Leg uit waarom een grotere beginsnelheid tot twee mogelijke
passagetijdstippen op $4,00\,\mathrm{m}$ kan leiden.
\end{question}

\end{oefening}


% ============================================================
% 13 MODEL EN LUCHTWEERSTAND
% ============================================================

\FVDsection{Wanneer faalt het vrije-valmodel?}

De vergelijkingen

\[
a_y=-g,
\qquad
v_y=v_0-gt,
\qquad
y=y_0+v_0t-\frac12gt^2
\]

zijn krachtig, maar niet universeel.

In lucht werkt naast de zwaartekracht ook luchtweerstand. Die hangt
onder andere af van vorm, oppervlakte, snelheid en eigenschappen van
de lucht.

Daardoor is de resulterende kracht niet meer constant en is de
versnelling niet noodzakelijk gelijk aan $-g$.


\FVDsubsection{Eindsnelheid}

Bij een vallend voorwerp kan de luchtweerstand toenemen wanneer de
snelheid groter wordt.

Wanneer de luchtweerstand uiteindelijk even groot is als de
zwaartekracht:

\[
F_{\mathrm{res}}=0.
\]

Dan:

\[
a=0
\]

en verandert de snelheid niet meer. Het voorwerp beweegt dan met een
\textbf{eindsnelheid}.

Dit gedrag kan niet worden beschreven door het eenvoudige
vrije-valmodel met constante versnelling $g$.


\begin{oefening}[Waarom is de voorspelling onrealistisch?]{\oefU\ESS}

Een model zonder luchtweerstand voorspelt voor een persoon die vanuit
rust $2000\,\mathrm{m}$ valt:

\[
|v|=\sqrt{2gh}.
\]

\begin{question}
Bereken de voorspelde snelheid.
\end{question}

\begin{question}
Zet het resultaat om naar $\mathrm{km/h}$.
\end{question}

\begin{question}
Waarom moet je voorzichtig zijn met de fysische interpretatie van
dit resultaat?
\end{question}

\begin{question}
Welke ontbrekende kracht maakt het model bij zo'n beweging
onvoldoende?
\end{question}

\end{oefening}


\begin{oefening}[Papier als modeltest]{\oefG}

Ontwerp een eenvoudig experiment met een blad papier waarmee je
aantoont dat vorm en luchtweerstand een belangrijke rol kunnen spelen.

Beschrijf:

\begin{question}
welke grootheid je verandert;
\end{question}

\begin{question}
welke grootheden je zoveel mogelijk constant houdt;
\end{question}

\begin{question}
welke waarneming je verwacht;
\end{question}

\begin{question}
waarom het resultaat geen tegenspraak vormt met het vrije-valmodel.
\end{question}

\end{oefening}


% ============================================================
% 14 DATA EN ICT
% ============================================================

\FVDsection{Een valbeweging onderzoeken met data}

Het leerplan vraagt niet alleen dat je formules toepast. Je moet ook
verbanden tussen positie, snelheid en versnelling kunnen analyseren.

Een videoanalyse is daarvoor bijzonder geschikt.

Film bijvoorbeeld een kleine bal die over een beperkte hoogte valt.
Bepaal met een videoanalysetool op opeenvolgende tijdstippen de
verticale positie.

Je krijgt meetpunten:

\[
(t_i,y_i).
\]

Voor vrije val verwachten we een kwadratisch model:

\[
y(t)=y_0+v_0t-\frac12gt^2.
\]

Bij loslaten uit rust:

\[
v_0\approx0.
\]

Uit het kwadratische model kunnen we door afleiden een lineaire
snelheidsfunctie en een constante versnellingsfunctie reconstrueren.


\begin{oefening}[Videoanalyse van een vallende bal]{\oefU\ESS}

Registreer of gebruik een gegeven video van een kleine bal die wordt
losgelaten.

\begin{question}
Kies een oorsprong en positieve richting en noteer die expliciet.
\end{question}

\begin{question}
Maak een tabel met $t$ en $y$.
\end{question}

\begin{question}
Maak een spreidingsdiagram van $y$ als functie van $t$.
\end{question}

\begin{question}
Bepaal met ICT een kwadratisch model
\[
y(t)=At^2+Bt+C.
\]
\end{question}

\begin{question}
Interpreteer $A$, $B$ en $C$ fysisch.
\end{question}

\begin{question}
Bepaal uit het model $v_y(t)$ en $a_y(t)$.
\end{question}

\begin{question}
Vergelijk de experimenteel gevonden versnelling met
$9,81\,\mathrm{m/s^2}$.
\end{question}

\begin{question}
Bespreek mogelijke oorzaken van een verschil.
\end{question}

\end{oefening}


% ============================================================
% 15 GEINTEGREERDE OEFENINGEN
% ============================================================

\FVDsection{Geïntegreerde oefeningen}

\begin{oefening}[Een voetbal verticaal omhoog]{\oefB\ESS}

Een voetbal wordt vanaf de grond verticaal omhoog getrapt met

\[
v_0=15,0\,\mathrm{m/s}.
\]

Verwaarloos luchtweerstand.

\begin{question}
Stel $y(t)$, $v_y(t)$ en $a_y(t)$ op.
\end{question}

\begin{question}
Bepaal wanneer de bal zijn hoogste punt bereikt.
\end{question}

\begin{question}
Bepaal de maximale hoogte.
\end{question}

\begin{question}
Bepaal wanneer de bal opnieuw de grond bereikt.
\end{question}

\begin{question}
Bepaal de snelheid vlak vóór het bereiken van de grond.
\end{question}

\begin{question}
Schets de drie bewegingsgrafieken.
\end{question}

\end{oefening}


\begin{oefening}[Een onbekende beginsnelheid]{\oefU\ESS}

Een bal wordt vanaf $2,00\,\mathrm{m}$ hoogte verticaal omhoog
geworpen. Hij bereikt een maximale hoogte van $7,00\,\mathrm{m}$.

\begin{question}
Bepaal de beginsnelheid.
\end{question}

\begin{question}
Bepaal de tijd tot het hoogste punt.
\end{question}

\begin{question}
Bepaal de totale tijd tot de bal de grond bereikt.
\end{question}

\begin{question}
Leg uit waarom je voor de laatste vraag de stijgtijd niet zomaar mag
verdubbelen.
\end{question}

\end{oefening}


\begin{oefening}[Grafiek zonder formule]{\oefU\ESS}

Van een verticale worp is alleen het volgende bekend:

\begin{itemize}
  \item $v_y(0)=19,62\,\mathrm{m/s}$;
  \item de $v(t)$-grafiek is een rechte met helling
        $-9,81\,\mathrm{m/s^2}$;
  \item $y(0)=1,00\,\mathrm{m}$.
\end{itemize}

\begin{question}
Bepaal grafisch of met de oppervlaktemethode de hoogtewinst tot het
hoogste punt.
\end{question}

\begin{question}
Bepaal de maximale hoogte.
\end{question}

\begin{question}
Controleer met een integraal.
\end{question}

\begin{question}
Stel daarna pas de plaatsfunctie op.
\end{question}

\end{oefening}


\begin{oefening}[Twee bollen]{\oefG}

Vanaf dezelfde hoogte wordt bol A losgelaten. Op hetzelfde tijdstip
wordt bol B verticaal naar beneden geworpen.

Beide bollen bewegen volgens het ideale vrije-valmodel.

\begin{question}
Vergelijk hun versnellingen.
\end{question}

\begin{question}
Stel voor beide bollen een snelheidsfunctie op.
\end{question}

\begin{question}
Onderzoek hoe het snelheidsverschil tussen beide bollen in de tijd
verandert.
\end{question}

\begin{question}
Onderzoek hoe hun onderlinge afstand in de tijd verandert.
\end{question}

\begin{question}
Formuleer een algemene conclusie en verklaar die fysisch.
\end{question}

\end{oefening}


\begin{oefening}[Foutenanalyse]{\oefU\ESS}

Een leerling lost een verticale worp op met omhoog positief:

\[
v_0=12,0\,\mathrm{m/s}.
\]

Hij schrijft:

\[
v(t)=12,0+9,81t
\]

en zegt:

\begin{quote}
``De bal gaat omhoog, dus $g$ moet positief zijn.''
\end{quote}

\begin{question}
Identificeer de fout.
\end{question}

\begin{question}
Schrijf de correcte snelheidsfunctie.
\end{question}

\begin{question}
Leg uit waarom de richting van de snelheid niet bepaalt welk teken de
versnelling heeft.
\end{question}

\end{oefening}


\begin{oefening}[Een model vergelijken met data]{\oefG}

Een sensor levert voor een vallend voorwerp meetgegevens op. Een
leerling vindt via regressie:

\[
y(t)=2,01-0,08t-4,72t^2.
\]

De positieve richting is omhoog.

\begin{question}
Bepaal de gemodelleerde beginsnelheid.
\end{question}

\begin{question}
Bepaal de gemodelleerde versnelling.
\end{question}

\begin{question}
Vergelijk deze versnelling met de theoretische waarde $-g$.
\end{question}

\begin{question}
Geef minstens drie mogelijke verklaringen voor het verschil.
\end{question}

\begin{question}
Beoordeel welke extra informatie je nodig hebt om te besluiten of
luchtweerstand de belangrijkste oorzaak is.
\end{question}

\end{oefening}


% ============================================================
% 16 SAMENVATTING
% ============================================================

\FVDsection{Wat hebben we opgebouwd?}

Een werkelijke valbeweging kan beïnvloed worden door verschillende
krachten. In het vrije-valmodel verwaarlozen we andere krachten dan
de zwaartekracht.

Dicht bij het aardoppervlak:

\[
g\approx9,81\,\mathrm{m/s^2}.
\]

Met omhoog positief:

\[
\boxed{a_y=-g.}
\]

Voor een verticale beweging met beginvoorwaarden $y_0$ en $v_0$:

\[
\boxed{
v_y(t)=v_0-gt
}
\]

en

\[
\boxed{
y(t)=y_0+v_0t-\frac12gt^2.
}
\]

Bij loslaten uit rust:

\[
v_0=0.
\]

Bij een verticale worp omhoog:

\[
v_0>0.
\]

Op het hoogste punt:

\[
\boxed{
v_y=0
\qquad\text{maar}\qquad
a_y=-g.
}
\]

Voor een worp omhoog:

\[
\boxed{
t_{\max}=\frac{v_0}{g}
}
\]

en:

\[
\boxed{
\Delta y_{\max}=\frac{v_0^2}{2g}.
}
\]

De relaties uit hoofdstuk 2 blijven volledig geldig:

\[
\boxed{
v_y(t)=y'(t),
\qquad
a_y(t)=v_y'(t)
}
\]

en:

\[
\boxed{
\Delta y=\int v_y(t)\,dt,
\qquad
\Delta v_y=\int a_y(t)\,dt.
}
\]

De georiënteerde oppervlakten onder de $v(t)$- en $a(t)$-grafieken
geven respectievelijk verplaatsing en snelheidsverschil.




% ============================================================
% 18 VOORUITBLIK
% ============================================================

\FVDsection{Vooruitblik}

Tot nu toe bewoog het voorwerp langs één verticale rechte.

Maar wat gebeurt er wanneer een bal horizontaal wordt weggeschoten?
De zwaartekracht versnelt hem verticaal, terwijl hij tegelijk
horizontaal beweegt.

Dan volstaat één coördinaat niet meer.

In het volgende hoofdstuk breiden we onze beschrijving uit naar
bewegingen in twee dimensies. We onderzoeken hoe onafhankelijke
bewegingscomponenten samen één baan vormen en gebruiken dat inzicht
om de horizontale worp te analyseren.

\end{document}
